\documentclass[11pt,a4paper]{article}

\usepackage[utf8]{inputenc}
\usepackage[T1]{fontenc}
\usepackage[french]{babel}
\usepackage{lmodern}
\usepackage{amsmath,amssymb,amsfonts}
\usepackage{geometry}
\usepackage{booktabs}
\usepackage{hyperref}
\usepackage{xcolor}
\usepackage{enumitem}
\usepackage{microtype}
\usepackage{tcolorbox}
\usepackage{array}

\geometry{margin=2.4cm}
\hypersetup{
  colorlinks=true,
  linkcolor=blue!60!black,
  urlcolor=blue!60!black,
  citecolor=blue!60!black,
  pdftitle={Consequences electromagnetiques de la conjecture UV/IR sur alpha},
  pdfauthor={Document de travail}
}

\setlength{\parskip}{0.55em}
\setlength{\parindent}{0pt}
\setlist[itemize]{leftmargin=1.4em}

\newcommand{\ee}{\mathrm e}
\newcommand{\cL}{\mathcal L}
\newcommand{\lp}{\ell_{\rm P}}
\newcommand{\RL}{R_{\Lambda}}
\newcommand{\me}{m_e}
\newcommand{\mP}{m_{\rm P}}
\newcommand{\lambdabar}{\bar\lambda_e}
\newcommand{\epsz}{\varepsilon_0}

\title{Consequences de la conjecture UV/IR sur $\alpha$\\
\large Charge, atome d'hydrogene, constantes electriques et cellule position--vitesse}
\author{Document de travail}
\date{22 mai 2026}

\begin{document}
\maketitle

\begin{abstract}
Ce document verifie les implications elementaires de la conjecture
\[
\alpha^{-1}=\ln\left(\frac{R_{\Lambda}}{10\pi\ell_{\rm P}}\right).
\]
En posant
\[
\mathcal L\equiv\ln\left(\frac{R_\Lambda}{10\pi\ell_{\rm P}}\right),
\]
la conjecture devient \(\alpha=1/\mathcal L\). Les grandeurs classiques de
l'electromagnetisme quantique se reecrivent alors en fonction du meme
logarithme UV/IR : charge elementaire, vitesse de Bohr, rayon de Bohr, energie
de Rydberg, rayon classique de l'electron, section de Thomson, resistance de
von Klitzing et rapport force electrique/force gravitationnelle. Le statut
reste conditionnel : ces relations sont des identites algebriques une fois
\(\alpha=1/\mathcal L\) pose, mais la preuve physique exige encore de deriver
le facteur \(10\pi\) depuis les modes de bord \(U(1)\).
\end{abstract}

\tableofcontents
\newpage

\section{Noyau conjectural}

On definit
\begin{equation}
\boxed{
\cL=\ln\left(\frac{\RL}{10\pi\lp}\right)
}
\end{equation}
avec
\begin{equation}
\RL=\sqrt{\frac{3}{\Lambda}},
\qquad
\lp=\sqrt{\frac{\hbar G}{c^3}}.
\end{equation}

La conjecture centrale est
\begin{equation}
\boxed{\alpha=\frac{1}{\cL}},
\qquad
\boxed{\alpha^{-1}=\cL}.
\end{equation}

\begin{tcolorbox}[colback=blue!3!white,colframe=blue!50!black,title=Statut]
Toutes les relations ci-dessous sont des consequences algebriques de la
substitution \(\alpha=1/\cL\). Elles ne constituent pas des preuves
independantes de la conjecture. Le verrou theorique demeure la derivation du
facteur \(10\pi\), ou du seuil \(S_0^{U(1)}=100\pi^3 k_B\), depuis la
microphysique de bord du secteur \(U(1)\).
\end{tcolorbox}

\section{Charge elementaire}

La constante de structure fine est
\begin{equation}
\alpha=\frac{e^2}{4\pi\epsz\hbar c}.
\end{equation}
Donc
\begin{equation}
e^2=4\pi\epsz\hbar c\,\alpha.
\end{equation}
Avec \(\alpha=1/\cL\),
\begin{equation}
\boxed{
e=\sqrt{\frac{4\pi\epsz\hbar c}{\cL}}.
}
\end{equation}

La charge de Planck vaut
\begin{equation}
q_{\rm P}=\sqrt{4\pi\epsz\hbar c}.
\end{equation}
Ainsi
\begin{equation}
\boxed{
e=\frac{q_{\rm P}}{\sqrt{\cL}}.
}
\end{equation}

Dans cette lecture, la charge elementaire est une charge de Planck affaiblie
par une racine logarithmique cosmologique.

\section{Force de Coulomb}

La force de Coulomb entre deux charges elementaires est
\begin{equation}
F_e=\frac{1}{4\pi\epsz}\frac{e^2}{r^2}.
\end{equation}
Comme \(e^2/(4\pi\epsz)=\alpha\hbar c\), on obtient
\begin{equation}
\boxed{
F_e=\frac{\hbar c}{\cL\,r^2}.
}
\end{equation}

\section{Vitesse de l'electron dans l'hydrogene}

Dans le modele de Bohr, la vitesse caracteristique de l'electron dans l'etat
fondamental de l'hydrogene est
\begin{equation}
v_e\simeq \alpha c.
\end{equation}
Donc
\begin{equation}
\boxed{
v_e=\frac{c}{\cL}.
}
\end{equation}

Cette relation doit etre lue comme une echelle non relativiste. La
spectroscopie de precision exige ensuite la masse reduite, l'equation de
Dirac, la structure fine, le Lamb shift et les corrections QED.

\section{Rayon de Bohr et longueur de Compton reduite}

Le rayon de Bohr s'ecrit
\begin{equation}
a_0=\frac{\hbar}{m_e c\alpha}.
\end{equation}
Avec \(\alpha=1/\cL\),
\begin{equation}
\boxed{
a_0=\frac{\hbar}{m_e c}\,\cL.
}
\end{equation}

La longueur de Compton reduite de l'electron est
\begin{equation}
\lambdabar=\frac{\hbar}{m_e c}.
\end{equation}
Donc
\begin{equation}
\boxed{
a_0=\lambdabar\,\cL.
}
\end{equation}

\section{Principe position--vitesse}

Dans l'atome d'hydrogene, au niveau d'echelle de Bohr,
\begin{equation}
a_0 v_e=
\left(\frac{\hbar}{m_e c}\cL\right)
\left(\frac{c}{\cL}\right).
\end{equation}
Donc
\begin{equation}
\boxed{
a_0 v_e=\frac{\hbar}{m_e}.
}
\end{equation}

Cela retrouve la forme position--vitesse du principe d'incertitude :
\begin{equation}
\Delta x\,\Delta v\sim\frac{\hbar}{m}.
\end{equation}

Dans le cadre GUP position--vitesse, la cellule effective s'ecrit
\begin{equation}
\boxed{
\Delta^3x\,\Delta^3v
\sim
\left(\frac{h}{m}\right)^3(1+\beta m^2v^2)^3.
}
\end{equation}

\section{Energie de Rydberg}

Dans la limite proton infiniment massif, l'energie fondamentale de l'hydrogene
est
\begin{equation}
E_1=-\frac12 m_e c^2\alpha^2.
\end{equation}
Avec \(\alpha=1/\cL\),
\begin{equation}
\boxed{
E_1=-\frac{m_e c^2}{2\cL^2}.
}
\end{equation}
L'energie de Rydberg correspondante est
\begin{equation}
\boxed{
E_R=\frac{m_e c^2}{2\cL^2}.
}
\end{equation}

Pour l'hydrogene reel, \(m_e\) doit etre remplacee par la masse reduite au
premier ordre.

\section{Rayon classique de l'electron}

Le rayon classique de l'electron est
\begin{equation}
r_e=\frac{e^2}{4\pi\epsz m_e c^2}.
\end{equation}
Comme \(e^2/(4\pi\epsz)=\alpha\hbar c\),
\begin{equation}
r_e=\alpha\frac{\hbar}{m_e c}.
\end{equation}
Donc
\begin{equation}
\boxed{
r_e=\frac{\lambdabar}{\cL}.
}
\end{equation}

On obtient la hierarchie
\begin{equation}
\boxed{
r_e=\frac{\lambdabar}{\cL},
\qquad
\lambdabar,
\qquad
a_0=\lambdabar\cL
}
\end{equation}
et donc
\begin{equation}
\boxed{
a_0 r_e=\lambdabar^2.
}
\end{equation}

\section{Section efficace de Thomson}

La section efficace de Thomson est
\begin{equation}
\sigma_T=\frac{8\pi}{3}r_e^2.
\end{equation}
En utilisant \(r_e=\lambdabar/\cL\),
\begin{equation}
\boxed{
\sigma_T=\frac{8\pi}{3}\frac{\lambdabar^2}{\cL^2}.
}
\end{equation}

\section{Champ critique de Schwinger}

Le champ critique de Schwinger est
\begin{equation}
E_S=\frac{m_e^2c^3}{e\hbar}.
\end{equation}
Le champ de Coulomb typique au rayon de Bohr verifie
\begin{equation}
E_{\rm Bohr}=\alpha^3 E_S.
\end{equation}
Donc
\begin{equation}
\boxed{
\frac{E_{\rm Bohr}}{E_S}=\frac{1}{\cL^3}.
}
\end{equation}

\section{Constantes electriques quantiques}

Le quantum de resistance de von Klitzing est
\begin{equation}
R_K=\frac{h}{e^2}.
\end{equation}
L'impedance du vide est
\begin{equation}
Z_0=\frac{1}{\epsz c}.
\end{equation}
On a
\begin{equation}
R_K=\frac{Z_0}{2\alpha}.
\end{equation}
Avec \(\alpha=1/\cL\),
\begin{equation}
\boxed{
R_K=\frac{Z_0}{2}\cL.
}
\end{equation}

Le quantum de conductance est
\begin{equation}
G_0=\frac{2e^2}{h}.
\end{equation}
Donc
\begin{equation}
\boxed{
G_0=\frac{4}{Z_0\cL}.
}
\end{equation}

Le quantum de flux supraconducteur est
\begin{equation}
\Phi_0=\frac{h}{2e}.
\end{equation}
Avec \(e=q_P/\sqrt{\cL}\),
\begin{equation}
\boxed{
\Phi_0=\frac{h}{2q_P}\sqrt{\cL}.
}
\end{equation}

\section{Rapport force electrique / force gravitationnelle}

Entre deux electrons,
\begin{equation}
F_G=\frac{Gm_e^2}{r^2}.
\end{equation}
Le rapport vaut
\begin{equation}
\frac{F_e}{F_G}
=
\frac{\alpha\hbar c}{Gm_e^2}.
\end{equation}
Avec \(\alpha=1/\cL\) et \(m_P=\sqrt{\hbar c/G}\),
\begin{equation}
\boxed{
\frac{F_e}{F_G}
=
\frac{m_P^2}{m_e^2}\frac{1}{\cL}.
}
\end{equation}

\section{Masse d'egalite electromagnetique--gravitationnelle}

On cherche une masse \(m_\star\) telle que
\begin{equation}
\frac{Gm_\star^2}{\hbar c}=\alpha.
\end{equation}
Donc
\begin{equation}
m_\star=m_P\sqrt{\alpha}.
\end{equation}
Avec \(\alpha=1/\cL\),
\begin{equation}
\boxed{
m_\star=\frac{m_P}{\sqrt{\cL}}.
}
\end{equation}

\section{Correction GUP pour l'electron atomique}

Dans la reformulation position--vitesse, le parametre de correction GUP est
\begin{equation}
\chi=\beta m^2v^2.
\end{equation}
Pour l'electron atomique,
\begin{equation}
v_e=\alpha c=\frac{c}{\cL}.
\end{equation}
En posant
\begin{equation}
\beta=\frac{\beta_0}{m_P^2c^2},
\end{equation}
on obtient
\begin{equation}
\boxed{
\chi_e=\beta_0\left(\frac{m_e}{m_P}\right)^2\frac{1}{\cL^2}.
}
\end{equation}

Numeriquement,
\begin{equation}
\frac{\chi_e}{\beta_0}\simeq 9.33\times10^{-50}.
\end{equation}
Pour \(\beta_0\sim1\), cette correction est invisible. En revanche, une valeur
gigantesque de \(\beta_0\), comme celle obtenue par identification directe
\(\rho_{\rm reg}=\rho_{\rm DE}\), serait exclue si elle etait universelle.

\section{Tableau synthetique}

\begin{center}
\renewcommand{\arraystretch}{1.35}
\begin{tabular}{>{\raggedright\arraybackslash}p{5.2cm} >{\raggedright\arraybackslash}p{8cm}}
\toprule
Grandeur & Reecriture avec \(\cL=\alpha^{-1}\) \\
\midrule
Charge elementaire & \(e=q_P/\sqrt{\cL}\) \\
Vitesse atomique de l'electron & \(v_e=c/\cL\) \\
Rayon de Bohr & \(a_0=\bar\lambda_e\cL\) \\
Rayon classique de l'electron & \(r_e=\bar\lambda_e/\cL\) \\
Produit fondamental & \(a_0r_e=\bar\lambda_e^2\) \\
Rydberg & \(E_R=m_ec^2/(2\cL^2)\) \\
Thomson & \(\sigma_T=(8\pi/3)\bar\lambda_e^2/\cL^2\) \\
Resistance de von Klitzing & \(R_K=(Z_0/2)\cL\) \\
Quantum de conductance & \(G_0=4/(Z_0\cL)\) \\
Rapport \(F_e/F_G\) entre electrons & \(F_e/F_G=m_P^2/(m_e^2\cL)\) \\
Masse d'egalite couplage gravitationnel/EM & \(m_\star=m_P/\sqrt{\cL}\) \\
Correction GUP atomique & \(\chi_e=\beta_0(m_e/m_P)^2/\cL^2\) \\
\bottomrule
\end{tabular}
\end{center}

\section{Interpretation}

Le coeur de toutes ces relations est
\begin{equation}
\boxed{
\cL=\ln\left(\frac{R_\Lambda}{10\pi\ell_P}\right)=\alpha^{-1}.
}
\end{equation}
Si cette conjecture est correcte, l'electromagnetisme atomique basse energie
devient une projection logarithmique du rapport Planck--de Sitter.

Mais il faut separer deux niveaux :
\begin{itemize}
\item le niveau algebrique, qui est verifie ;
\item le niveau physique, qui exige encore une derivation du seuil
\(S_0^{U(1)}=100\pi^3k_B\).
\end{itemize}

\section{Verrou theorique}

Dans la formulation entropique, le facteur \(10\pi\) correspond a
\begin{equation}
S_0^{U(1)}=100\pi^3k_B.
\end{equation}
Une piste est
\begin{equation}
S_0^{U(1)}=4\pi\left[2\pi\left(2+\frac12\right)\right]^2 k_B,
\end{equation}
ou
\begin{equation}
2=\text{deux polarisations du photon},
\qquad
\frac12=\text{mode de bord }U(1)\text{ renormalise}.
\end{equation}

Le programme de preuve consiste donc a deriver rigoureusement
\begin{equation}
\boxed{
g_{\rm edge}^{U(1)}=\frac12
}
\end{equation}
comme capacite holographique renormalisee du mode de bord Maxwell.

\section{Conclusion}

Les formules de cette note sont coherentes et utiles comme cartographie des
consequences de la conjecture \(\alpha^{-1}=\ln(R_\Lambda/(10\pi\ell_P))\).
Elles montrent que charge, echelle atomique, diffusion Thomson, resistance
quantique et hierarchie electromagnetisme-gravitation peuvent etre reecrites a
partir du meme logarithme UV/IR.

La conclusion scientifique stricte est cependant prudente : ces reecritures ne
valident pas seules la conjecture. Elles indiquent ou chercher des tests et des
contraintes, en particulier la non-variation temporelle de \(\alpha\), les
limites atomiques sur \(\beta_0\), et le calcul de bord Maxwell.

\begin{thebibliography}{9}

\bibitem{CODATA2022}
NIST/CODATA, \emph{CODATA recommended values of the fundamental physical
constants: 2022}.

\bibitem{Planck2018}
Planck Collaboration, \emph{Planck 2018 results. VI. Cosmological parameters},
arXiv:1807.06209.

\bibitem{DonnellyWall2015}
W. Donnelly and A. C. Wall, \emph{Geometric entropy and edge modes of the
electromagnetic field}, arXiv:1506.05792.

\end{thebibliography}

\end{document}
